\scanpage{94}
\mspage{11}
Mais avant, un commentaire. Quand on donne $\underline I$, rev. cubique étale de $S$, et une section $t$ de $\Sigma_{\underline I}^*$, alors la donnée d'une courbe elliptique relative sur $S$, correspondante à ces données $\unsure{\underline I|\underline\Pi|\underline J},t$, équivaut à celle d'un module inv. $T$ sur $S$, et d'un isom $(5)$
\[
T^{\otimes2}\simeq T_{\Sigma,t} \tag*{(*)}
\]
Quand on a dispose déjà d'un tel $T$, soit $T_0$, alors la catégorie des $T$ \inserted{sur $S$} munis de $(*)$ équivaut à celle des $\mu_{2\,S}$-torseurs $\underline\varepsilon$ sur $S$, moyennant l'équivalence
\[
\deleted{\unsure{\underline T}}\underline\varepsilon\longmapsto T=T_0\mathbin{\wedge}_{\mu_{2\,S}}\underline\varepsilon \tag*{(55)}
\]
Ici, on a envie de prendre
\[
T_0=\underline O_\Sigma(1)(t)^{\unsure{\sharp|\xi|\natural}}=\underline O_\Sigma(1)(t)\mathbin{\wedge}_{\mu_4}Q_{0S} \tag*{(56)}
\]
mais \deleted{\unsure{p}}on a un iso. can.
\[
T_0^{\otimes2}\simeq\underline O_\Sigma(2)(t)\mathbin{\wedge}_{\mu_2}\mu_4^*
\simeq T_{\Sigma,t}\mathbin{\wedge}_{\mu_2}(\mu_4^\times\mathbin{\wedge}\underline\omega_{\underline I})
\]
et pour avoir un iso. ``raisonnable'' $T_0^{\otimes2}\simeq T_{\Sigma,t}$, il faut se donner une \deleted{section} iso
\[
\alpha\unsure{\in|:}\deleted{\unsure{\sharp}}\mu_4^\times\simeq\underline\omega_{\underline I}
\quad\text{i.e. une section de }\mu_4^\times\mathbin{\wedge}\underline\omega_{\underline I}
\]
(ce qui est une des données indiquées par un épinglage de Legendre d'une courbe elliptique $/S$ correspondant à $\underline I$ \ldots). Une fois $\alpha$ fixé, on
% archive sheet number: 93
